\section{Industry and Laboratory Experience}

\entry{Ph.D.\ Extern}{May 2025 -- Aug.\ 2025}{IBM Research}{}
\begin{itemize}
  \item Designed a portable collective communication library for AI workloads
    (all-gather, all-reduce, reduce-scatter), building on and contributing to the
    open-source UCCL library for moving large tensors between devices
  \item Combined UCCL with Charm++ so that collectives tolerate accelerator failures and
    adapt to any cluster's hardware and network topology; demonstrated on a distributed
    tensor multiplication written in Charm++
\end{itemize}
\entrygap

\entry{Ph.D.\ Intern}{May 2023 -- Aug.\ 2023}{Pacific Northwest National Laboratory}{}
\begin{itemize}
  \item Developed a minimal-communication asynchronous algorithm for distributed
    reinforcement learning training on GPUs, built with RLlib on the Ray API, and
    applied it to speed up training for the influence maximization problem
\end{itemize}
\entrygap

\entry{Control Systems Engineering Co-op}{Jan.\ 2021 -- June 2021}{Nuvera Fuel Cells}{}
\begin{itemize}
  \item Developed Simulink control software for a hydrogen fuel cell engine, raising
    power output from 45 to 60~kW and lowering the required starting temperature from
    0\,\textdegree C to $-30$\,\textdegree C
  \cvonly{\item Ported fuel cell stack partial-derivative calculations from MATLAB to C,
    and ran lab tests analyzing performance data to tune the control code}
\end{itemize}
\entrygap

\entry{GPU Software Engineering Co-op}{Jan.\ 2020 -- June 2020}{Advanced Micro Devices (AMD)}{}
\begin{itemize}
  \item Designed a tool that instruments GPU shader programs during compilation to record
    performance metrics at run time
  \cvonly{\item Worked with the shader compiler team to define project goals and
    specifications, and surveyed similar tools from other companies}
\end{itemize}
